\section{The characterization}\label{sec:introduction}

The cotype--cotype problem asks whether finite Rademacher cotype
of a Banach space and its dual forces \(K\)-convexity. Without an
approximation hypothesis the assertion is false. We prove the characterization under the
ordinary approximation property.

All spaces in this paper are real, and \(X^*\) denotes the continuous
dual of \(X\). Recall that \(X\) has the \emph{approximation property}
(AP) if, for every compact set \(M\subset X\) and every \(\delta>0\),
there is a bounded finite-rank linear map \(S:X\to X\) such that
\[
\sup_{x\in M}\|Sx-x\|<\delta.
\]
The norms of these maps need not have a common bound. The
\emph{bounded approximation property} (BAP) requires, in addition,
a constant \(\lambda<\infty\) such that all the maps needed for these
approximations can be chosen with \(\|S\|\le\lambda\).

For \(2\le q<\infty\), let \(C_q(Y)\) be the least constant such that
\[
\left(\sum_{i=1}^s\|y_i\|^q\right)^{1/q}
\le C_q(Y)
\left(\E_\eps\left\|\sum_{i=1}^s\eps_i y_i\right\|^2\right)^{1/2}
\]
for every finite family in \(Y\), where the \(\eps_i\) are independent
uniform signs. The space \(Y\) has \emph{finite Rademacher cotype} if
\(C_q(Y)<\infty\) for some finite \(q\).

The \emph{Rademacher projection} on \(X\)-valued functions on a sign
cube retains their first-order coefficients:
\[
(\Rad_s f)(\eps)=\sum_{i=1}^s\eps_i\,\E_\eta[\eta_i f(\eta)].
\]
Write
\[
K(X)=\sup_{s\ge1}
\|\Rad_s:L_2(\{-1,1\}^s;X)\to L_2(\{-1,1\}^s;X)\|.
\]
Here \(L_2\) uses uniform probability measure. By definition, \(X\) is
\emph{\(K\)-convex} when \(K(X)<\infty\).

\begin{theorem}\label{thm:main}
Let \(X\) be a nonzero real Banach space with the ordinary approximation
property. Then \(X\) is \(K\)-convex if and only if
\[
C_q(X)<\infty
\quad\hbox{and}\quad
C_r(X^*)<\infty
\qquad\text{for some }q,r\in[2,\infty).
\]
The exponents \(q\) and \(r\) may differ. The forward implication holds
for every real Banach space, without an approximation hypothesis.
\end{theorem}

A space \(X\) has \emph{Rademacher type \(p\)}, \(1<p\le2\), if for some \(D<\infty\),
\[
\left(\E_\eps\left\|\sum_i\eps_i x_i\right\|^2\right)^{1/2}
\le D\left(\sum_i\|x_i\|^p\right)^{1/p}
\]
for all finite families. We say that \(X\) contains the \(\ell_1^s\)'s
uniformly if the isomorphism distortions of such copies inside \(X\)
are bounded independently of \(s\).

The classical theorems of Maurey--Pisier and Pisier
\cite{MP76,Pis82K} identify \(K\)-convexity with nontrivial type:
\begin{equation}\label{eq:classical-characterization}
\begin{split}
K(X)<\infty
&\ \Longleftrightarrow\
X\text{ has Rademacher type }p\text{ for some }p>1\\
&\ \Longleftrightarrow\
X\text{ does not contain the }\ell_1^s\text{'s uniformly}.
\end{split}
\end{equation}
These equivalences require no approximation hypothesis; precise
formulations are \cite[Theorem~4.40(i)]{Pis11IHP} and
\cite[Theorem~1.5]{Pis80Kseminar}.
The finite-cube definition above agrees with the infinite-product
definition in \cite[Definition~1.1]{Pis80Kseminar}, by density of
functions depending on finitely many coordinates.

The characterization is uniform when the cotype data are fixed. Its
finite-dimensional form also supplies a class of convex bodies for which
ordinary metric-entropy duality has dimension-free constants.
For an origin-symmetric convex body \(B\subset\R^n\), write \(X_B\)
for \(\R^n\) with unit ball \(B\), and \(B^\circ\) for its polar.
Here a convex body is compact with nonempty interior, and \(N(S,T)\)
is the least number of translates of \(T\) needed to cover \(S\).

\begin{corollary}[Uniform cotype bounds and entropy duality]
\label{cor:uniform-cotype-entropy}
Fix \(q,r\in[2,\infty)\) and \(C,D\in[1,\infty)\). There is a
finite constant \(\kappa=\kappa(q,r,C,D)\) such that every nonzero
real Banach space \(X\) with AP and
\[
C_q(X)\le C,\qquad C_r(X^*)\le D
\]
satisfies \(K(X)\le\kappa\).
Consequently there are \(a,b\ge1\), depending only on \(q,r,C,D\),
such that for every \(n\ge1\), all origin-symmetric convex bodies
\(K,L\subset\R^n\), and every \(t>0\),
\begin{equation}\label{eq:cotype-entropy-duality}
\frac1b\log N(L^\circ,atK^\circ)
\le \log N(K,tL)
\le b\log N(L^\circ,a^{-1}tK^\circ),
\end{equation}
provided either \(X_K\) or \(X_L\) satisfies the displayed primal
and dual cotype bounds.
\end{corollary}

The constants are existential; no bound on the AP approximants is
required. The entropy conclusion uses the bounded-\(K\) duality theorem
of Artstein--Milman--Szarek--Tomczak-Jaegermann
\cite[Theorem~5]{AMSTJ2004}, in the covering-number form of
\cite[Theorem~1.1 and Section~2]{Milman2007}.
We prove the corollary at the end of Section~\ref{sec:approximation}.

\subsection{History and scope}

Type and cotype express the geometry of a Banach space through the
size of random signed sums. Their relation to uniform finite-dimensional
subspaces was developed by Maurey and Pisier \cite{MP76}.
Elementary duality passes a type estimate to a cotype estimate on the
dual, but recovering type from cotype requires further information.
Maurey and Pisier introduced \(K\)-convexity as the setting for this
duality question; Pisier's holomorphic-semigroup theorem subsequently
identified it with nontrivial type
\cite{MP76,Pis80Kseminar,Pis82K}.
The question here asks whether finite cotype on both sides, together
with approximation, already puts the space in that setting.

Approximation assumptions have played a substantial role in these questions.
Enflo's counterexample showed that a separable reflexive Banach space can
fail AP \cite{Enf73}. A Schauder basis gives uniformly bounded finite-rank
approximations, but the converse implications in the hierarchy
\[
\text{Schauder basis}\ \Longrightarrow\ \text{BAP}\ \Longrightarrow\ \text{AP}
\]
fail: Figiel and Johnson separated AP from BAP \cite[p.~199]{FJ73}, and
Szarek constructed a superreflexive space with BAP but no Schauder basis
\cite[Theorem~1.1]{Sza87}.
Thus an argument using ordinary AP must accommodate approximants whose
norms have no common bound.

For cotype two on both sides, Pisier established the stronger conclusion
that an AP space is isomorphic to a Hilbert space
\cite[Corollaries~2.9 and~2.16]{Pis80}.
His subsequent work proved \(K\)-convexity under a restriction on the
cotype exponents, first with a basis or a dense increasing family of
uniformly complemented finite-dimensional subspaces, and then with BAP
\cite[Theorem~1, Remarks~2 and~9]{Pis81}.
He asked whether the exponent restriction could be removed
\cite[Remark~2(iii)]{Pis81}.
Szarek and Tomczak-Jaegermann discuss this question as the
\emph{cotype--cotype conjecture} in connection with uniformly complemented
finite-dimensional subspaces \cite{STJ09}.
Gupta, Misra, and Ray state the BAP formulation in
\cite[Conjecture~1.5]{GMR25}.
Theorem~\ref{thm:main} proves the characterization for all finite cotype
exponents under ordinary AP.

The assertion with no approximation hypothesis is false.
Pisier's construction embeds every cotype-two space into a space whose
primal and dual both have cotype two \cite[Theorem~3.2]{Pis83}.
Starting with \(\ell_1\) therefore gives a space which is not
\(K\)-convex; these examples fail AP \cite[Remark~3.6]{Pis83}.

Another application concerns a uniform tensor estimate. Given a sequence
\(\mathcal E=(E_n)\) with \(\dim E_n=n\), one asks whether every map
\(A:E_m\to E_n^*\), tensored with the identity of \(F\), is bounded
from the injective tensor norm to the projective tensor norm by
one constant times \(\|A\|\), independently of \(m,n,A\).
This is the Grothendieck-pair condition, stated formally in
Section~\ref{sec:grothendieck}. Gupta--Misra--Ray prove that it forces
finite cotype of \(F\) and \(F^*\), and record Pisier's conjecture
that for infinite-dimensional BAP spaces \(F\) the estimate then also
holds with \(\ell_2\) in place of \(F\)
\cite[Definition~1.2, Theorem~1.3 and Conjecture~1.4]{GMR25}.
Their deduction through complemented Euclidean subspaces
\cite[Corollary~4.8]{GMR25}, combined with Theorem~\ref{thm:main},
gives this conclusion under ordinary AP.
This uniform operator estimate differs from Grothendieck's original
injective/projective tensor-product conjecture, disproved in
\cite{Pis83}, and from the GT-pairs of \cite[Section~6]{Pis12}.

Rademacher and Walsh expansions connect probabilistic moment estimates
with Banach-space geometry. Kahane's inequality compares moments of
vector-valued Rademacher sums \cite[Theorem~8.1]{Pis11IHP};
Borell extended moment comparison to Banach-valued Walsh polynomials
of bounded degree \cite{Borell79}. The moment lifts below combine the
Rademacher comparison with quotient duality and Hahn--Banach. The scalar coefficient
lemma uses an additional feature of Fourier indices: their rank over
\(\F_2\), which determines the number of independent Walsh factors.

\subsection{The mechanism of the proof}

The proof has two parts. First, we establish a uniform decay estimate
for the normalized Walsh transforms of finite-rank operators
\(T:G\to F\): apply \(T\) to a vector family and then multiply by the
matrix \((2^{-n}(-1)^{u\cdot v})_{u,v\in\F_2^n}\), with the dot product
computed modulo two and probability-normalized \(L_2\) norms.
Cotype estimates for \(G^*\) and \(F\) at one sufficiently large, fixed
tuple length suffice. The decay estimate, relative to \(\|T\|\), is independent of the
operator rank and of the spaces; Theorem~\ref{thm:uniform-decay} states the
precise result.

A fixed finite-rank operator already has a Walsh bound obtained by
factoring through a Hilbert space, but its constant depends on the
operator. To remove that dependence, we show that a slowly decreasing
Walsh norm forces another substantial norm at nearly twice the
dimension. A moment lift represents a tuple \((x_i)\) by a bounded
sign-indexed family \((P_\eps)\) with
\(\E_\eps\eps_iP_\eps=ax_i\), where \(a>0\) is controlled by cotype.
Such lifts and finite-group averaging turn the slow decrease into
bounded scalar functions with prescribed row and column moments.
Their Fourier expansion is indexed by binary matrices and must contain
an index of rank at least two over \(\F_2\). Such an index produces
two independent Walsh factors.
Iteration would then contradict the Hilbert-space bound for the same
fixed operator. This proves uniform decay.

Second, we transfer that decay to the identity on \(X\) using only AP.
For each finite-dimensional \(E\subset X\), we construct a compact map
\(J:G\to X\) of uniformly bounded norm and a contraction \(\iota:E\to G\)
with \(J\iota x=x\). The dual of \(G\) has a controlled moment estimate
(Lemma~\ref{lem:compact-domain}). The construction uses moment lifts in
\(X\), rather than assuming that the quotient \(E^*\) inherits the
cotype of \(X^*\). AP gives finite-rank composites \(S_jJ\) converging
to \(J\) in operator norm. Uniform Walsh decay passes to \(J\), and
then to every finite Walsh input in \(X\). Uniform copies of
\(\ell_1^s\) would contradict that decay.

Section~\ref{sec:preliminaries} introduces Walsh norms and the moment lifts.
Section~\ref{sec:coefficient} proves the scalar coefficient lemma, and
Section~\ref{sec:walsh} turns it into uniform operator decay.
Section~\ref{sec:approximation} constructs the compact domains and proves
Theorem~\ref{thm:main} and Corollary~\ref{cor:uniform-cotype-entropy}.
Section~\ref{sec:grothendieck} gives the tensor-norm consequence.
